\documentclass[11pt]{article}
\usepackage{amsmath,amssymb,amsthm,mathrsfs}
\usepackage[margin=1in]{geometry}
\usepackage{hyperref}

\theoremstyle{plain}
\newtheorem{theorem}{Theorem}
\newtheorem{lemma}[theorem]{Lemma}
\newtheorem{proposition}[theorem]{Proposition}
\newtheorem{corollary}[theorem]{Corollary}
\theoremstyle{definition}
\newtheorem{definition}[theorem]{Definition}
\newtheorem{remark}[theorem]{Remark}
\newtheorem{conjecture}[theorem]{Conjecture}

\newcommand{\A}{\mathscr{A}}
\newcommand{\QH}{\mathrm{QH}^*}
\newcommand{\Gr}{\mathrm{Gr}}
\newcommand{\ZZ}{\mathbb{Z}}
\DeclareMathOperator{\hgt}{ht}

\title{The height of a ribbon is the orientation of a heap:\\
the Hall--Littlewood $t$-family inside the affine\\ nil-Temperley--Lieb algebra}
\author{Clio}
\date{23 September 2026 (Day 201, PROVE cycle 2)}

\begin{document}
\maketitle

\begin{center}\fbox{\begin{minipage}{0.92\textwidth}
\small
\textbf{Correspondence block (PROTOCOL \S2.3).}\quad
\textbf{Author:} Clio.\quad
\textbf{Date of this build:} 24 September 2026.\quad
\textbf{Recipient:} Rick.\\[2pt]
\textbf{Title:} The height of a ribbon is the orientation of a heap.\quad
\textbf{Commit:} \texttt{clio-vega/proofs@0a01bb9}.\\[4pt]
The mathematics of this note is that of commit \texttt{0a01bb9} (23 September 2026,
PROVE cycle 2) and \emph{has not changed since}: no statement, proof, definition,
numerical claim or range is altered. Three non-mathematical changes were made on
24 September 2026, all in commit \texttt{a7b218e}, and they are listed by kind because
``nothing has changed'' is a claim about every other line and survives nothing:
(i) this block was added;
(ii) \verb|\newcommand{\ht}| \emph{silently failed} --- \verb|\ht| is a \TeX{} primitive,
so all six uses were parsed as box-dimension lookups and rendered as nothing, in a note
whose subject is ribbon height; it is now \verb|\DeclareMathOperator{\hgt}{ht}|;
(iii) the \verb|]| inside the optional argument of Theorem~\ref{thm:A} closed it early,
truncating its title. \texttt{pdflatex} now completes with zero errors, where the
23~September build had twenty-four.
\\[4pt]
\textbf{What is proved and what is not.} Theorem~\ref{thm:A} and
Theorem~\ref{thm:C} are \emph{proved}. Theorem~\ref{thm:B} is \emph{computed, not
proved} --- see gap (G1) --- and Proposition~\ref{prop:MN} is conditional on it.
\\[4pt]
\textbf{Novelty is unchecked.} No priority claim is made anywhere in this note. Theorem~A
has \emph{not} been checked against Korff--Vasilev \texttt{2606.06352},
Benkart--Meinel \texttt{1505.02544}, or the affine nil-Temperley--Lieb literature
(OpenAlex returns sixty works, none of which was in my sources index before
23 September 2026). Three sentences are forbidden in any draft and are not asserted
here: that the operator-level cylindric--integrable bridge is new (Korff--Stroppel 2010
own it); that this is the first combinatorial proof of Postnikov's symmetry; and that
Postnikov's proof is unwritten (it is on the page).
\end{minipage}}\end{center}

\bigskip

\begin{abstract}
Let $\A_n$ be Postnikov's affine nil-Temperley--Lieb algebra, acting on
$\bigoplus_k\QH(\Gr_{kn})$ by adding boxes to cylindric diagonals.
Postnikov's elements $\mathbf e_r$ and $\mathbf h_r$ are built from
\emph{two} special orderings of each generator set.  We show that the
intermediate orderings are not noise: for an interval $J\subset\ZZ/n\ZZ$
of length $e$, the $2^{e-1}$ distinct monomials of $\A_n$ with support
$J$ are in canonical bijection with the $2^{e-1}$ ribbons of size $e$,
and the number of edges of $J$ oriented ``upward'' in the monomial's
commutation class is exactly the \emph{height} (leg length) of the
ribbon.  Consequently the whole height-graded ribbon-addition family
$R_e(t)=\sum_h t^h N_e^{(h)}$ lives in $\A_n[t]$ as an orientation sum,
$\mathbf h_e$ and $\mathbf e_e$ contain its $t=0$ and $t=\infty$
specialisations as their connected parts, and the functional equation
$\Omega R_e(t)\Omega=t^{e-1}R_e(1/t)$ is the Grassmannian duality
$\mathbf e\leftrightarrow\mathbf h$.  We then locate the two quantum
features of Postnikov's Lemma (the vanishing $\mathbf e_i\mathbf h_j=0$
for $i+j>n$, and the correction term at order $t^n$): both are the single
relation $\mathbf e_k\mathbf h_{n-k}=q$, and on the ribbon side the
mechanism producing them is the ribbon that wraps all the way round the
cylinder, $e=n$.  Finally we record strong computational evidence that
$t=-1$ is the \emph{unique} point of the $t$-line at which $R_e(t)$
commutes with the cylindric strip-adders --- the cylindric counterpart of
the classical statement that $R_e(t)$ is a multiplication operator
exactly at $t=-1$ --- and that Newton's identity holds in $\A_n$ with
$p_e:=R_e(-1)$ for $r\le n-1$, failing at $r=n$ by exactly the wrap term.
All computations are exact and symbolic; ranges are stated with each
claim.
\end{abstract}

\section{The question, and what is proved}

\subsection*{Background}

DREAM 2026-09-23 c1 framed question Q237 as: Postnikov asserts
(\texttt{math/0205165v2} ll.1460--1464) an unwritten combinatorial proof
of the $S_\ell$-symmetry of cylindric skew Schur functions via commuting
\emph{strip operators}; are those my ribbon-adding vertex operators
$R_e(t)$?  It recorded as a blocker: \emph{``Postnikov's sketch is not on
the page, so nothing here can be quoted.''}  That blocker is false: the
operators are written down in $\A_n$ (ll.2837--2879) and their
commutation is a lemma with an explicit defect term (ll.2882--2893).  The
sharpened target, set at WAKE 2026-09-23 c2, was:

\begin{quote}
\textbf{Target.} Exhibit an index identification and a specialisation of
$t$ under which the commutator identity satisfied by $R_e(t)$ specialises
to Postnikov's \texttt{eq:e-h-AN}, \emph{including} its truncation
$\mathbf e_i\mathbf h_j=0$ for $i+j>n$; or exhibit an index-set
obstruction that kills the identification.
\end{quote}

\subsection*{Answers}

\begin{itemize}
\item \textbf{Index set (\S\ref{sec:heaps}).}  The two index sets are the
  same set, and neither is a quotient of the other.  Postnikov's $\sim$
  (l.2795) is the cylinder's shift; his own l.2806 states that
  $\kappa\mapsto\mathbf a_\kappa$ is a bijection between cylindric shapes
  mod $\sim$ and nonvanishing monomials of $\A_n$.  What is added here is
  the \emph{statistic}: the commutation class of the monomial records the
  ribbon's height.  Theorem~\ref{thm:A}.
\item \textbf{Parameter (\S\ref{sec:t}).}  There is no tension between
  ``nil'' and $t=-1$.  The nil relations constrain which \emph{states} a
  monomial acts on, not the height; the full $t$-line sits inside
  $\A_n[t]$.  Postnikov's $\mathbf h_e,\mathbf e_e$ sit at $t=0$ and
  $t=\infty$; and $t=-1$ --- Rick's fermion-bilinear point --- is where
  $R_e(t)$ becomes a multiplication operator and \texttt{eq:e-h-AN}
  becomes Newton's identity.  Theorem~\ref{thm:B}.
\item \textbf{Truncation (\S\ref{sec:trunc}).}  There \emph{is} a
  mechanism on the ribbon side, and it is the wrap-around ribbon of size
  exactly $n$, which Postnikov's index set excludes by fiat (his $I$ must
  be a \emph{proper} subset of $\ZZ/n\ZZ$).  The correction term equals
  $(-1)^{n-k}q\cdot\mathrm{Id}$ on $\QH(\Gr_{kn})$ --- the single quantum
  relation $\mathbf e_k\mathbf h_{n-k}=q$.  Theorem~\ref{thm:C}.
\end{itemize}

So the identification is not killed; the obstruction is on the other
side, and it is an obstruction to the \emph{semi-infinite wedge}, whose
index set $\ZZ$ has no wrap.

\section{Setup and conventions}

\subsection{The algebra and its action}

Following Postnikov \cite{Pos} (\texttt{math/0205165v2}; all line numbers
below refer to the \LaTeX{} e-print I hold, at
\texttt{sources/math0205165-postnikov/postnikov-affine-v2.tex}, and were
read at source), for $n\ge 2$ the \emph{affine nil-Temperley--Lieb
algebra} $\A_n$ (ll.2647--2657) is the associative $\ZZ$-algebra with
generators $a_i$, $i\in\ZZ/n\ZZ$, and relations
\begin{equation}\label{eq:rel}
a_ia_i=a_ia_{i+1}a_i=a_{i+1}a_ia_{i+1}=0,\qquad
a_ia_j=a_ja_i\ \ (i-j\not\equiv\pm1).
\end{equation}

A partition $\lambda\in P_{kn}$ is encoded by its $01$-word
$\omega(\lambda)\in\{0,1\}^n$ with $k$ ones (l.302): $0=$ right step,
$1=$ up step of the boundary path.  Equivalently, writing
$\beta_j=\lambda_j+k-j$ for $1\le j\le k$, the \emph{beads} of $\lambda$
occupy the sites $\beta_1>\dots>\beta_k$ in $\{0,\dots,n-1\}$.  We index
sites by $\ZZ/n\ZZ$ with $0$-based labels throughout, so Postnikov's
$a_1,\dots,a_n$ are our $a_0,\dots,a_{n-1}$.  The action (ll.2670--2698)
is
\[
a_i=E_{i,i+1}\ (0\le i\le n-2),\qquad a_{n-1}=q\cdot E_{n-1,0},
\]
where $E_{ij}$ moves a bead from site $i$ to site $j$ when $i$ is
occupied and $j$ is empty, and kills the basis vector otherwise.  Thus:

\begin{quote}
\emph{$a_i$ moves one bead forward from site $i$ to site $i+1$, requiring
$i$ occupied and $i+1$ empty, and contributes a factor $q$ exactly when
it crosses the seam $n-1\to 0$.}
\end{quote}

\begin{remark}[a load-bearing convention]\label{rem:conv}
Postnikov defines (ll.2839--2848) $\prod^{\circlearrowright}_{i\in I}a_i$
as the product of the $a_i$, $i\in I$, ``taken in an order such that if
$i,i+1\in I$ then $a_{i+1}$ goes before $a_i$'', and displays
(l.2862) $\mathbf e_2=a_2a_1+a_3a_2+\cdots$; he also states (l.2877) that
$\mathbf e_r$ enumerates cylindric \emph{vertical} $r$-strips.  Under the
standard right-to-left composition convention these two are
incompatible: with $a_1$ applied first, $a_2a_1$ sends $\lambda=(2)$ to
$(2,2)$ in $P_{2,5}$, a horizontal strip.  They are consistent if the
displayed products are read in \emph{application} order (leftmost factor
applied first), equivalently if ``goes before'' means ``is applied
before''.  We adopt that reading; in our own (right-to-left) notation
Postnikov's $\prod^{\circlearrowright}$ over a run $[i,i+e-1]$ is written
$a_i a_{i+1}\cdots a_{i+e-1}$.  This is recorded because the statistic in
Theorem~\ref{thm:A} is \emph{precisely} a statement about application
order, so the convention is load-bearing rather than cosmetic.
\end{remark}

\subsection{Heaps: monomials with a given support}

Call $I\subseteq\ZZ/n\ZZ$ \emph{proper} if $I\ne\ZZ/n\ZZ$.  For proper
$I$, the subgraph of the $n$-cycle induced on $I$ is a disjoint union of
paths; write $\mathrm{Ed}(I)$ for its edge set, so
$|\mathrm{Ed}(I)|=|I|-(\text{number of maximal runs of }I)$.

\begin{definition}\label{def:aIeps}
For proper $I$ and $\varepsilon\in\{0,1\}^{\mathrm{Ed}(I)}$, let
$a_{I,\varepsilon}\in\A_n$ be the product of $\{a_i:i\in I\}$, each once,
in any linear order such that for every edge $\{s,s+1\}\in\mathrm{Ed}(I)$,
\[
\varepsilon_{s}=0\ \Longrightarrow\ a_{s}\text{ is applied before }a_{s+1},
\qquad
\varepsilon_{s}=1\ \Longrightarrow\ a_{s+1}\text{ is applied before }a_{s}.
\]
(Edges are indexed by their smaller endpoint in each run.)
\end{definition}

$a_{I,\varepsilon}$ is well defined: two linear orders satisfying the
same $\varepsilon$ differ by transpositions of non-adjacent generators,
which commute by \eqref{eq:rel}.  Conversely two linear orders that
disagree on some edge give different elements of the trace monoid.  Thus
the monomials of $\A_n$ with support $I$ (each generator used once) are
exactly the $2^{|\mathrm{Ed}(I)|}$ elements $a_{I,\varepsilon}$.  In
Postnikov's notation,
\begin{equation}\label{eq:extremes}
\textstyle\prod^{\circlearrowleft}_{i\in I}a_i=a_{I,\mathbf 0},
\qquad
\prod^{\circlearrowright}_{i\in I}a_i=a_{I,\mathbf 1}.
\end{equation}

\subsection{Ribbons and heights}

For a skew shape $\theta$ that is a rim hook (connected, no $2\times2$),
write $\hgt(\theta)=(\text{number of rows of }\theta)-1$ for its leg
length.  We use the classical abacus dictionary in the following form,
which we verified independently (\S\ref{sec:verif}, Check 1):

\begin{lemma}\label{lem:abacus}
Let $\lambda$ have beads $S$, let $j\in S$ and $j+e\notin S$ with
$j+e\le n-1$, and let $\mu$ have beads $(S\setminus\{j\})\cup\{j+e\}$.
Then $\mu/\lambda$ is a rim hook of size $e$, and
\[
\hgt(\mu/\lambda)=\#\{s: 1\le s\le e-1,\ j+s\in S\},
\]
the number of beads the moving bead jumps over.  The arm is
$e-1-\hgt$.
\end{lemma}

Accordingly we define, for $1\le e$, the \emph{cylindric height-graded
$e$-ribbon adder} directly on beads:
\begin{equation}\label{eq:ribbondef}
\mathcal R_e(t)\,\sigma_S \;=\;
\sum_{\substack{j\in S\\ j+e\notin S \bmod n}}
t^{\,\#\{1\le s\le e-1\,:\,j+s\in S \bmod n\}}\;
q^{\,[\,n-1\in\{j,\dots,j+e-1\}\,]}\;
\sigma_{(S\setminus\{j\})\cup\{j+e\}},
\end{equation}
all indices mod $n$.  For $e\le n-1$ this is the cylindric avatar of the
operator $R_e(t)=\sum_{h=0}^{e-1}t^hN_e^{(h)}$ of
\texttt{topics/hall-littlewood-vertex-operator-dictionary.md}: when no
wrap occurs it is, by Lemma~\ref{lem:abacus}, literally $\lambda\mapsto
\sum_\mu t^{\hgt(\mu/\lambda)}\mu$ over $e$-ribbon additions.

\section{Theorem A: the height is the orientation}\label{sec:heaps}

\begin{lemma}[Interval monomials are ribbons]\label{lem:run}
Let $1\le e\le n-1$, let $J=[i,i+e-1]\subset\ZZ/n\ZZ$ be an interval, and
let $\varepsilon\in\{0,1\}^{e-1}$ (edges indexed by
$s=1,\dots,e-1$, edge $s$ joining $i+s-1$ and $i+s$).  Then, on a bead
configuration $S$, the monomial $a_{J,\varepsilon}$ acts as follows:
$a_{J,\varepsilon}\sigma_S\ne 0$ if and only if
\begin{equation}\label{eq:pattern}
i\in S,\qquad i+e\notin S,\qquad
\text{and }\ i+s\in S\iff \varepsilon_s=1\ \ (1\le s\le e-1),
\end{equation}
and in that case
$a_{J,\varepsilon}\sigma_S=q^{[\,n-1\in J\,]}\,\sigma_{(S\setminus\{i\})\cup\{i+e\}}$.
In particular $|\varepsilon|$ is the number of beads jumped over, i.e.\
the height of the ribbon added.
\end{lemma}

\begin{proof}
Each generator $a_s$, $s\in J$, is applied exactly once.  Among these
generators, the occupancy of a site $u$ is changed only by $a_{u-1}$
(which fills $u$) and by $a_u$ (which empties $u$); and $a_{u-1}\in J$
iff $u-1\in J$.

\emph{Necessity.}  Site $i$: only $a_i\in J$ affects it (since
$a_{i-1}\notin J$), and $a_i$ requires $i$ occupied at its firing time;
as nothing fills $i$ beforehand, $i\in S$.  Site $i+e$: only
$a_{i+e-1}\in J$ affects it, and it requires $i+e$ empty at its firing
time; as nothing empties $i+e$ beforehand, $i+e\notin S$.  Now fix
$1\le s\le e-1$ and put $u=i+s$; the only generators in $J$ touching $u$
are $a_{u-1}$ and $a_u$, so the occupancy of $u$ immediately before the
first of these fires equals its initial occupancy.  If
$\varepsilon_s=1$ then $a_u$ fires first, and it requires $u$ occupied,
so $u\in S$.  If $\varepsilon_s=0$ then $a_{u-1}$ fires first, and it
requires $u$ empty, so $u\notin S$.  This is \eqref{eq:pattern}.

\emph{Sufficiency.}  Assume \eqref{eq:pattern} and take any linear order
compatible with $\varepsilon$.  We check each $a_s$, $s\in J$, is legal
when its turn comes.

\emph{Source site $s$.}  If $s=i$: $i\in S$ and no generator of $J$ fills
or empties $i$ before $a_i$, so $i$ is occupied.  If $s=i+r$ with
$r\ge1$: when $\varepsilon_r=1$, $a_s$ precedes $a_{s-1}$ and
$s\in S$ by \eqref{eq:pattern}, so $s$ is still occupied; when
$\varepsilon_r=0$, $a_{s-1}$ precedes $a_s$ and $s\notin S$, so $a_{s-1}$
has already filled $s$.

\emph{Target site $s+1$.}  If $s=i+e-1$: $i+e\notin S$ and no generator
of $J$ fills $i+e$ before $a_{i+e-1}$, so it is empty.  If $s=i+r$ with
$r\le e-2$: when $\varepsilon_{r+1}=0$, $a_s$ precedes $a_{s+1}$ and
$s+1\notin S$, so $s+1$ is still empty; when $\varepsilon_{r+1}=1$,
$a_{s+1}$ precedes $a_s$ and $s+1\in S$, but $a_{s+1}$ has already
emptied it.

So every firing is legal, for \emph{any} compatible linear order.  Each
$a_s$ removes a bead from $s$ and puts one at $s+1$; summing over
$s\in J$ the intermediate sites telescope and the net change of the bead
set is $-i,+(i+e)$.  The scalar collected is $q$ for each $s\in J$ with
$s=n-1$, i.e.\ $q^{[\,n-1\in J\,]}$, and $|J|=e\le n-1$ so this exponent
is $0$ or $1$.  Finally $|\varepsilon|=\#\{s:i+s\in S,\ 1\le s\le e-1\}$
by \eqref{eq:pattern}.
\end{proof}

\begin{theorem}[{The $t$-family lies in $\A_n[t]$}]\label{thm:A}
Let $1\le e\le n-1$ and set
\begin{equation}\label{eq:Rdef}
\boxed{\;
R_e(t)\;:=\;\sum_{i\in\ZZ/n\ZZ}\ \sum_{\varepsilon\in\{0,1\}^{e-1}}
t^{\,|\varepsilon|}\; a_{[i,\,i+e-1],\,\varepsilon}\ \in\ \A_n[t].\;}
\end{equation}
Then $R_e(t)$ acts on $\bigoplus_k\QH(\Gr_{kn})$ exactly as the cylindric
height-graded $e$-ribbon adder $\mathcal R_e(t)$ of
\eqref{eq:ribbondef}.  Moreover, for fixed $i$ the $2^{e-1}$ summands act
on pairwise disjoint sets of basis vectors, so the grading by
$|\varepsilon|$ is the grading by height, term by term.
\end{theorem}

\begin{proof}
Fix $S$ and $i$ with $i\in S$, $i+e\notin S$.  The occupancy pattern of
$S$ on the interior sites $i+1,\dots,i+e-1$ determines a unique
$\varepsilon\in\{0,1\}^{e-1}$ via \eqref{eq:pattern}, and by
Lemma~\ref{lem:run} exactly that summand is nonzero on $\sigma_S$; it
contributes $t^{|\varepsilon|}q^{[n-1\in J]}\sigma_{(S\setminus\{i\})\cup\{i+e\}}$,
which is the $j=i$ term of \eqref{eq:ribbondef}.  If $i\notin S$ or
$i+e\in S$ then every summand at that $i$ vanishes, matching
\eqref{eq:ribbondef}.  Summing over $i$ gives the claim, and disjointness
of supports is \eqref{eq:pattern}.
\end{proof}

\begin{remark}
Theorem~\ref{thm:A} is compatible with, and sharpens, Postnikov's own
l.2806: \emph{``the map $\kappa\mapsto\mathbf a_\kappa$ gives a
one-to-one correspondence between cylindric shapes (modulo the
`$\sim$'-equivalence) and nonvanishing monomials in the algebra
$\A_n$''}, where $\sim$ (ll.2795--2803) is generated by independent
south-east shifts of connected components.  Restricted to
\emph{connected} shapes on $e$ consecutive diagonals --- i.e.\ to
$e$-ribbons --- that correspondence is exactly our
$\varepsilon\leftrightarrow a_{J,\varepsilon}$, because an $e$-ribbon is
determined by its word of $e-1$ steps in $\{\text{right},\text{down}\}$.
What is new is the identification of the \emph{statistic}: the down-steps
of the ribbon are the ``up''-oriented edges of the heap, hence
$\hgt=|\varepsilon|$, hence the entire $t$-line is an orientation
generating function inside one algebra.
\end{remark}

\begin{corollary}[Postnikov's $\mathbf e$ and $\mathbf h$ are the two ends
of the $t$-line]\label{cor:ends}
With $\mathbf e_r,\mathbf h_r$ as in \cite[ll.2853--2857]{Pos},
\[
R_e(0)=\sum_{i}\textstyle\prod^{\circlearrowleft}_{J=[i,i+e-1]}a
\ =\ (\text{connected part of }\mathbf h_e),
\qquad
[t^{e-1}]R_e(t)=\sum_i\textstyle\prod^{\circlearrowright}_{J}a
\ =\ (\text{connected part of }\mathbf e_e).
\]
Equivalently: a cylindric horizontal $e$-strip is a disjoint union of
non-adjacent height-$0$ ribbons and a cylindric vertical $e$-strip is a
disjoint union of non-adjacent height-maximal ribbons, and $\mathbf h_e$,
$\mathbf e_e$ are obtained from $R_{\bullet}(0)$, $R_\bullet(\infty)$ by
allowing several non-adjacent runs.
\end{corollary}

\begin{proof}
Immediate from \eqref{eq:extremes} and Lemma~\ref{lem:run}: the runs of
a proper subset $I$ are non-adjacent, hence their monomials commute and
the bead moves they perform are disjoint; $\varepsilon=\mathbf 0$ gives
height $0$ on each run and $\varepsilon=\mathbf 1$ gives height
$\ell_a-1$ on a run of length $\ell_a$.
\end{proof}

\begin{corollary}[The duality $\mathbf e\leftrightarrow\mathbf h$ is
$t\mapsto 1/t$]\label{cor:duality}
Let $\rho$ be the Grassmannian duality $P_{kn}\to P_{n-k,n}$,
$\lambda\mapsto\lambda'$, which on beads is
$\rho(S)=\{\,n-1-j \bmod n\ :\ j\notin S\,\}$.  Then
\[
\rho\, \mathcal R_e(t)\,\rho^{-1}\;=\;t^{\,e-1}\,\mathcal R_e(1/t).
\]
\end{corollary}
\noindent
This is the cylindric form of the functional equation $\Omega
R_e(t)\Omega=t^{e-1}R_e(1/t)$ (PROVE 2026-08-30).  It is verified
symbolically for $3\le n\le 7$, all $1\le k\le n-1$ (\S\ref{sec:verif},
Check 5); a proof is immediate from Lemma~\ref{lem:abacus}
($\rho$ exchanges arm and leg) but is not written out here.
Corollaries~\ref{cor:ends} and \ref{cor:duality} together say: the
$\mathbf e/\mathbf h$ duality of $\A_n$ \emph{is} the involution
$t\mapsto1/t$ of the ribbon $t$-line, exchanging its two endpoints.

\section{Theorem B: the parameter is $t=-1$}\label{sec:t}

Corollary~\ref{cor:ends} places $\mathbf h$ and $\mathbf e$ at $t=0$ and
$t=\infty$.  The interior of the $t$-line carries the structure that
makes Postnikov's Lemma an \emph{identity} rather than a definition.

\begin{theorem}[verified, ranges as stated]\label{thm:B}
Let $2\le k\le n-2$.
\begin{enumerate}
\item[(i)] \emph{(Rigidity.)}  $t=-1$ is the unique value of $t$ for
  which $[\,\mathcal R_e(t),\mathbf h_j\,]=0$ for all $1\le e\le n-1$ and
  all $1\le j\le n-1$.  Verified for $4\le n\le 6$.  (For $k=1$ and
  $k=n-1$ the commutator vanishes identically in $t$; the space is too
  small to see a constraint, and those cases carry no evidence either
  way.)
\item[(ii)] \emph{(Newton.)}  Writing $p_e:=\mathcal R_e(-1)$, the
  operators $p_1,\dots,p_{n-1}$ commute with one another and with
  $\mathbf h_1,\dots,\mathbf h_{n-1}$, and
  \[
  r\,\mathbf h_r \;=\; \sum_{e=1}^{r} p_e\,\mathbf h_{r-e}
  \qquad (1\le r\le n-1).
  \]
  Verified for $3\le n\le 7$ and all $1\le k\le n-1$, exactly and
  symbolically in $\ZZ[q]$.
\item[(iii)] \emph{(The defect.)}  At $r=n$ the identity fails, with
  \[
  \sum_{e=1}^{n-1}p_e\,\mathbf h_{n-e}\;=\;(-1)^{k-1}(n-k)\,q\cdot\mathrm{Id}
  \qquad\text{on }\QH(\Gr_{kn}).
  \]
  Verified for $3\le n\le 7$, all $k$.
\end{enumerate}
\end{theorem}

\begin{proposition}[Consequence: the quantum Murnaghan--Nakayama rule]\label{prop:MN}
Assume Theorem~\ref{thm:B}(ii).  Then for $1\le e\le n-1$,
$\mathcal R_e(-1)$ is quantum multiplication by $p_e$ on
$\QH(\Gr_{kn})$:
\[
\mathcal R_e(-1)\,\sigma_\lambda \;=\; p_e * \sigma_\lambda .
\]
\end{proposition}

\begin{proof}
By Postnikov's quantum Pieri corollary (\texttt{cor:pieri-AN},
ll.2935--2947) $\mathbf h_j$ acts as quantum multiplication by
$h_j=\sigma_j$.  Let $\varphi:\Lambda[q]\to\QH(\Gr_{kn})$ be the ring
homomorphism determined by $e_i\mapsto\sigma_{1^i}$ (zero for $i>k$);
this is well defined because $\Lambda=\ZZ[e_1,e_2,\dots]$ is free, and by
the presentation \texttt{eq:QH} (ll.500--507) it is surjective and sends
$h_j\mapsto\sigma_j$ for $1\le j\le n-1$.  Newton's identity
$r\,h_r=\sum_{e\le r}p_eh_{r-e}$ holds in $\Lambda$, hence its
$\varphi$-image holds in $\QH(\Gr_{kn})$.  Now induct on $e$.  For $e=1$,
Theorem~\ref{thm:B}(ii) at $r=1$ gives
$\mathcal R_1(-1)=\mathbf h_1=M_{p_1}$.  For $e=r\le n-1$, subtracting
the two Newton relations,
\[
\mathcal R_r(-1)\;=\;r\,\mathbf h_r-\sum_{e=1}^{r-1}\mathcal R_e(-1)\mathbf h_{r-e}
\;=\;M_{r h_r}-\sum_{e=1}^{r-1}M_{p_e}M_{h_{r-e}}
\;=\;M_{r h_r-\sum_{e<r}p_eh_{r-e}}\;=\;M_{p_r}. \qedhere
\]
\end{proof}

\begin{remark}[why $t=-1$ and not $t=0$]
The brief anticipated a tension: $\A_n$ is \emph{nil} ($a_i^2=0$), which
in most such families is a $t\to0$ degeneration, so $t=-1$ and ``nil''
might be incompatible.  Lemma~\ref{lem:run} shows the tension is not
there.  The nil relations constrain the \emph{support} of a monomial ---
which occupancy patterns \eqref{eq:pattern} it can act on --- and not its
height; heights $0,\dots,e-1$ all occur, in the $2^{e-1}$ different
commutation classes with the same support $J$.  What is true is that each
individual monomial is supported on a single height; the $t$-grading is
recovered only after summing over the commutation classes, which is
exactly \eqref{eq:Rdef}.
\end{remark}

\section{Theorem C: where the truncation comes from}\label{sec:trunc}

Postnikov's Lemma (\texttt{lem:e-h-commute-in-AN}, ll.2882--2893) states
that $\mathbf e_i,\mathbf h_j$ commute pairwise, that
$\mathbf e_i\mathbf h_j=0$ for $i+j>n$, and that
\begin{equation}\label{eq:ehAN}
\Bigl(1+\sum_{i=1}^{n-1}\mathbf e_it^i\Bigr)
\Bigl(1+\sum_{j=1}^{n-1}\mathbf h_j(-t)^j\Bigr)
=1+\Bigl(\sum_{m=1}^{n-1}(-1)^{n-m}\mathbf e_m\mathbf h_{n-m}\Bigr)t^n .
\end{equation}

\begin{theorem}\label{thm:C}
On $\QH(\Gr_{kn})$, $1\le k\le n-1$:
\begin{enumerate}
\item[(i)] the correction term of \eqref{eq:ehAN} is the scalar operator
  \[
  \sum_{m=1}^{n-1}(-1)^{n-m}\mathbf e_m\mathbf h_{n-m}
  \;=\;(-1)^{n-k}\,q\cdot\mathrm{Id};
  \]
\item[(ii)] it is nothing but the single quantum relation
  $\mathbf e_k\mathbf h_{n-k}=q$, all other terms of the sum being zero;
\item[(iii)] on the ribbon side, the images under $\varphi$ of the two
  ``missing'' degree-$n$ objects are
  \[
  \varphi(h_n)=(-1)^{k-1}q,\qquad \varphi(p_n)=(-1)^{k-1}k\,q ,
  \]
  and $\varphi(p_n)$ is exactly the value of the \emph{full-wrap} bead
  move: each of the $k$ beads travels $n$ steps round the cylinder,
  jumping over the other $k-1$ beads and crossing the seam once, giving
  $k\cdot t^{k-1}q\big|_{t=-1}$.
\end{enumerate}
\end{theorem}

\begin{proof}
(i) and (ii).  By Postnikov's presentation \texttt{eq:QH} (ll.500--507),
$\QH(\Gr_{kn})=\ZZ[q,e_1,\dots,e_k,h_1,\dots,h_{n-k}]/\langle
e(t)h(-t)=1+(-1)^{n-k}qt^n\rangle$, with $e_i=0$ for $i>k$ and $h_j=0$
for $j>n-k$.  The coefficient of $t^n$ of the defining relation reads
$\sum_{m=0}^{n}(-1)^{n-m}e_mh_{n-m}=(-1)^{n-k}q$.  In the sum, $e_m=0$
unless $m\le k$ and $h_{n-m}=0$ unless $n-m\le n-k$, i.e.\ $m\ge k$; so
the only surviving term is $m=k$, giving $(-1)^{n-k}e_kh_{n-k}=(-1)^{n-k}q$,
i.e.\ $e_kh_{n-k}=q$.  Since $1\le k\le n-1$, the index $m=k$ lies in the
range $1\le m\le n-1$ of \eqref{eq:ehAN}, and all other terms vanish.
By \texttt{cor:pieri-AN} the operators $\mathbf e_m,\mathbf h_j$ are the
corresponding quantum multiplications, so the operator identity follows.

(iii).  With $\varphi$ as in Proposition~\ref{prop:MN}: in $\Lambda$ we
have $h(-t)=e(t)^{-1}$, whereas in $\QH(\Gr_{kn})$
$h_{\QH}(-t)=e(t)^{-1}\bigl(1+(-1)^{n-k}qt^n\bigr)$.  Hence, writing
$u=-t$,
\[
\varphi\bigl(h(u)\bigr)=h_{\QH}(u)\bigl(1-(-1)^{k}q\,u^{n}+O(u^{2n})\bigr),
\]
and taking the coefficient of $u^n$, using $h_{\QH,n}=0$ and
$h_{\QH,0}=1$, gives $\varphi(h_n)=(-1)^{k-1}q$.  Newton at $r=n$ in
$\Lambda$, pushed through $\varphi$, then gives
\[
\varphi(p_n)=n\,\varphi(h_n)-\sum_{e=1}^{n-1}\varphi(p_e)\varphi(h_{n-e})
= n(-1)^{k-1}q-(-1)^{k-1}(n-k)q=(-1)^{k-1}k\,q,
\]
where the middle sum is evaluated by Proposition~\ref{prop:MN} together
with Theorem~\ref{thm:B}(iii).  The stated bead count gives the same
number: for each of the $k$ beads $j$, the move $j\mapsto j+n=j$ is legal
(the site is vacated by the bead itself), it passes the other $k-1$ beads
and crosses the seam once, contributing $t^{k-1}q$; at $t=-1$ the total
is $k(-1)^{k-1}q$.
\end{proof}

\begin{corollary}[the answer to the truncation question]
The vanishing $\mathbf e_i\mathbf h_j=0$ for $i+j>n$ and the correction
term of \eqref{eq:ehAN} are \textbf{not} features of the $t$-family.
They are features of the \emph{index set}.  Postnikov's index set is
$\{\text{proper subsets of }\ZZ/n\ZZ\}$; the ribbon family's index set
is $\{\text{cylindric ribbons}\}$, which contains the wrap-around ribbon
of size $n$ excluded by ``proper''.  On the semi-infinite wedge, whose
index set is $\ZZ$, there is no wrap and hence no truncation; the
mechanism that produces one is precisely the surjection
$\ZZ\twoheadrightarrow\ZZ/n\ZZ$.
\end{corollary}

\begin{remark}[one relation, four faces]
The relation $e_kh_{n-k}=q$ --- the only difference between $H^*(\Gr_{kn})$
and $\QH(\Gr_{kn})$ --- appears in this paper in four guises:
as the correction term $(-1)^{n-k}q$ of \eqref{eq:ehAN}
(Theorem~\ref{thm:C}(i));
as the failure of Newton's identity at $r=n$, of size $(-1)^{k-1}(n-k)q$
(Theorem~\ref{thm:B}(iii));
as $\varphi(h_n)=(-1)^{k-1}q$;
and as $\varphi(p_n)=(-1)^{k-1}kq$, the full-wrap ribbon
(Theorem~\ref{thm:C}(iii)).  The arithmetic
$n(-1)^{k-1}q=(-1)^{k-1}(n-k)q+(-1)^{k-1}kq$ that ties the last three
together is the whole quantum content, read off the cylinder as
``$n$ steps $=$ $(n-k)$ gaps $+$ $k$ beads''.
\end{remark}

\section{Computational verification}\label{sec:verif}

All checks are exact and symbolic (SymPy over $\ZZ[q,t]$), in
\texttt{proofs/code-q237/} (\texttt{anTL.py}, \texttt{ribbon\_check.py}).
The model implements $\A_n$ from the presentation
\eqref{eq:rel} \emph{via its action}, i.e.\ $a_i$ is the bead move, and
the relations \eqref{eq:rel} are consequences rather than inputs.  Two
independent routes to each object are compared.

\begin{enumerate}
\item \textbf{Abacus dictionary (Lemma~\ref{lem:abacus}).}  Non-cyclic
bead moves $j\to j+e$ versus the Young-diagram skew shape, checked for
size, connectedness, absence of $2\times2$, $\text{rows}-1=h$ and
$\text{cols}-1=e-1-h$.  All $n\le 9$, all $k$, all $S$, all legal moves:
\textbf{0 violations}.
\item \textbf{Theorem~\ref{thm:A}.}  Orientation sum \eqref{eq:Rdef}
versus the independent bead definition \eqref{eq:ribbondef}.
$2\le n\le 9$, all $0\le k\le n$, all $1\le e\le n-1$:
\textbf{0 mismatches}.
\item \textbf{Disjoint supports / nonvanishing.}  For $e=4$,
$n\in\{6,7\}$, $k\in\{2,3,4\}$: the $2^{e-1}=8$ monomials per interval
have pairwise disjoint state-supports (\textbf{0 overlaps}); those that
vanish identically are exactly the ones violating $|\varepsilon|+1\le k$
or $e-|\varepsilon|\le n-k$, as \eqref{eq:pattern} predicts.
\item \textbf{Postnikov's Lemma, reproduced.}  For $3\le n\le 7$, all
$k$: the coefficients of $t^1,\dots,t^{n-1}$ in the left side of
\eqref{eq:ehAN} vanish; $\mathbf e_i\mathbf h_j=0$ for $i+j>n$; the
$t^n$ coefficient equals Postnikov's stated right-hand side; and that
scalar is $(-1)^{n-k}q$.
\item \textbf{Corollary~\ref{cor:duality}.}  $\rho$ with
$\rho(S)=\{n-1-j:j\notin S\}$ intertwines $\mathcal R_e(t)$ on
$(n,k)$ with $t^{e-1}\mathcal R_e(1/t)$ on $(n,n-k)$, for
$3\le n\le7$, all $k$, $e=\min(3,n-1)$.
\item \textbf{Theorem~\ref{thm:B}.}  Commutation
$[\mathcal R_e(t),\mathbf h_j]$ solved for $t$ entrywise: the common root
set is $\{-1\}$ for $4\le n\le6$, $2\le k\le n-2$, and unconstrained for
$k\in\{1,n-1\}$.  Newton verified for $3\le n\le7$, all $k$, all
$r\le n-1$; defect at $r=n$ equals $(-1)^{k-1}(n-k)q$ in every case.
\item \textbf{Full-wrap value.}  The $e=n$ bead move evaluated at $t=-1$
equals $(-1)^{k-1}kq$ for $3\le n\le7$, all $k$.
\end{enumerate}

\section{Gaps}

These are stated as propositions, not as descriptions, so that each is
falsifiable on its own.

\begin{enumerate}
\item[(G1)] \textbf{Theorem~\ref{thm:B} is verified, not proved.}  Both
(i) rigidity at $t=-1$ and (ii) Newton's identity in $\A_n$ are open.
Proposition~\ref{prop:MN} derives the quantum Murnaghan--Nakayama rule
\emph{from} (ii); it does not prove (ii).  Equivalently one may prove
$[\mathcal R_e(-1),\mathbf h_j]=0$ directly, since with quantum Giambelli
(\texttt{eq:Giambelli}, ll.534--540) an operator commuting with all
$\mathbf h_j$ is multiplication by its value on $\sigma_\varnothing$.
I expect a sign-reversing involution on pairs
(cylindric horizontal strip, signed ribbon); I did not find it today.
\item[(G2)] \textbf{The $q=0$ half of (G1) is not the hard half.}  At
$q=0$, $\mathcal R_e(-1)=M_{p_e}$ on $H^*(\Gr_{kn})$ follows from the
classical statement $R_e(-1)=M_{p_e}$ on $\Lambda$ (PROVE 2026-09-03 c2,
Q75) because multiplication by $p_e$ preserves the ideal spanned by
Schur functions outside the box.  All the content of (G1) is in the
$q^1$ terms, i.e.\ in the wrapping ribbons.  I have \emph{not} checked
that the $q=0$ reduction argument is airtight; it is written here as a
plan, not as a proof.
\item[(G3)] \textbf{Rigidity is untested at the boundary.}  For $k=1$ and
$k=n-1$ the commutator $[\mathcal R_e(t),\mathbf h_j]$ vanishes for all
$t$.  I have \emph{not} determined whether this is a rank degeneracy of
small spaces or a genuine statement about $\Gr_{1,n}=\mathbb P^{n-1}$;
until that is decided, the phrase ``$t=-1$ is unique'' must carry the
hypothesis $2\le k\le n-2$.
\item[(G4)] \textbf{The convention of Remark~\ref{rem:conv} is a reading,
not a quotation.}  Postnikov's displayed $\mathbf e_2$ (l.2862) and his
combinatorial description (l.2877) are incompatible under right-to-left
composition; I resolved this in favour of the combinatorial description,
because that is what \texttt{cor:pieri-AN} and the rest of the section
use.  A reader holding a different build of the paper should re-check
l.2862 before relying on \eqref{eq:extremes}.  All line numbers here
refer to the e-print I hold; they are internal locators, not published
theorem numbers.
\item[(G5)] \textbf{Corollary~\ref{cor:duality} is verified and sketched,
not proved.}  The sketch (``$\rho$ exchanges arm and leg'') is one line
and I believe it; it is not written out.
\item[(G6)] \textbf{Nothing here touches the $S_\ell$-symmetry itself.}
Theorem~\ref{thm:A} identifies the operators; it does not give the
commutativity argument that Postnikov's ll.1460--1464 sketch, nor
does it bear on the M-convexity programme (WZZ Problem 5.2), which
concerns supports rather than operators.
\end{enumerate}

\section*{Novelty statement}

\textbf{The operator-level bridge between cylindric combinatorics and
integrable systems is not new, and no sentence here may claim it.}
Korff--Stroppel \cite{KS} (l.272) explicitly \emph{``reprov[e] a result of
Postnikov''} by realising $\A_n$ on a fermionic Fock space --- the
generators $a_i$ of this paper are, in their language, ``one fermion
hopping in clockwise direction on the circle'' (l.2474).  Korff--Vasilev,
\emph{Equivariant Quantum Cohomology of Grassmannians via the Clifford
algebra} (\texttt{arXiv:2606.06352}, May 2026), go further: \S
\texttt{sec:cylloops} (l.1483) puts Postnikov's cylindric loops in
periodic Maya diagrams --- the bead model of this paper --- and App.~B
(l.1734) gives the equivariant analogue of $\A_n$ as
$\varphi(u_i)=\psi^*_{i+1}\psi_i+y_i\psi^*_i\psi_i$.  Benkart--Meinel,
\emph{The center of the affine nil-Temperley--Lieb algebra}
(\texttt{arXiv:1505.02544}), study exactly this algebra through ``a
faithful representation of it on fermionic particle configurations''.
An OpenAlex query for ``affine nil-Temperley--Lieb'' returns sixty works,
\emph{none} of which was in my sources index before this morning.

I therefore make no priority claim of any kind.  What I claim is that
the following three statements were derived in this session and are, to
my knowledge, not written down in the sources I hold: the identification
of the ribbon \emph{height} with the orientation statistic of the heap
(Theorem~\ref{thm:A}); the consequent realisation of the whole $t$-line
$R_e(t)$ --- not just its two endpoints --- inside $\A_n[t]$; and the
location of \texttt{eq:e-h-AN}'s truncation as the $e=n$ wrap-around
ribbon (Theorem~\ref{thm:C}).  I have \textbf{not} checked any of the
three against \texttt{2606.06352}, \texttt{1505.02544}, or the sixty-work
nil-Temperley--Lieb literature; that check is owed, and until it is done
each of the three should be read as ``derived here'', not ``new''.

In particular the following sentences are forbidden in any draft leaving
this container: ``the operator-level cylindric--integrable bridge is
new''; ``the first combinatorial proof of Postnikov's symmetry'';
``Postnikov's proof is unwritten'' (the sketch at ll.1460--1464 is
unwritten; the algebra, the lemma and the identity are written, at
ll.2837--2893).

\begin{thebibliography}{9}
\bibitem{Pos} A.~Postnikov, \emph{Affine approach to quantum Schubert
calculus}, \texttt{arXiv:math/0205165v2}.  Cited by internal
\LaTeX{} labels and line numbers of the e-print source at
\texttt{sources/math0205165-postnikov/postnikov-affine-v2.tex},
read at source on 2026-09-23.
\bibitem{KS} C.~Korff and C.~Stroppel, \emph{The
$\widehat{\mathfrak{sl}}(n)_k$-WZNW fusion ring: a combinatorial
construction and a realisation as quotient of quantum cohomology},
Adv.\ Math.\ 225 (2010) 200--268, \texttt{arXiv:0909.2347}.  Line
numbers refer to
\texttt{sources/0909.2347-korff-stroppel/src/revisedversion4.tex}.
\end{thebibliography}

\end{document}
